\chapter{Completeness of the normalisation}\label{sec:lex}

Conditions (3)--(4) of Definition~\ref{def:phi} break symmetry: they demand a
normalised neighbourhood for the fixed vertex and lexicographic minimality of
the orbit word under $101$ relabelling generators.  This chapter proves that
these conditions lose no graphs: every $\sigma_0$-invariant
$\srg(99,14,1,2)$ has an isomorphic copy whose orbit word satisfies all of
them, and consequently $\Phi$ is satisfiable whenever such a graph exists.
The argument has three layers: transport of strong regularity along graph
isomorphisms, a minimisation over the orbit of the word under the generator
action (by well-ordering of $\mathbb{N}$), and the satisfiability of the
clauses that render the lexicographic constraints.

%---------------------------------------------------------------------
\section{Transport of strong regularity}

\begin{lemma}[Common neighbours across an isomorphism]\label{lem:lex-iso-common-neighbors}
A graph isomorphism $e \colon G \to H$ restricts, for every pair of vertices
$v,w$ of $G$, to a bijection between the common neighbours of $v,w$ in $G$
and the common neighbours of $e(v),e(w)$ in $H$.
\lean{ConwayO7.isoCommonNeighbors}\leanok
\uses{def:srg}
\end{lemma}



%---------------------------------------------------------------------
\section{The block-relabelling subgroup}

\begin{definition}[Relabelling generators]\label{def:relabel}
Let $N \le S_{99}$ be generated by: the block transposition exchanging
$7$-blocks $0,1$; the transpositions of adjacent blocks $i,i{+}1$ for
$2 \le i \le 12$; the rotations of each block; and the five multiplier
permutations $j \mapsto aj \bmod 7$ ($2 \le a \le 6$) applied simultaneously
to all blocks.  Each $\rho \in N$ normalises $\langle\sigma_0\rangle$, hence
acts on the orbit space and thus on words $\word{693}$ by
$(\rho\cdot\omega)(o) = \omega(\rho^{-1}o)$.  The encoding uses $101$ fixed
generators of this action.
\lean{ConwayO7.Encoder.lexPerms, ConwayO7.Encoder.orbitPerm}\leanok
\uses{lem:orbits}
\end{definition}

\begin{definition}[Block relabellings]\label{def:lex-cycle-relabel}
Write $B_i = \{1+7i,\dots,7+7i\}$ ($0 \le i \le 13$) for the fourteen
$7$-blocks rotated by $\sigma_0$.  A permutation $\tau \in S_{14}$ of the
block indices induces a permutation of the $99$ vertices fixing $0$ and
sending the $j$-th vertex of block $i$ to the $j$-th vertex of block
$\tau(i)$: block indices are permuted, positions within blocks are kept.
\lean{ConwayO7.Encoder.cycleRelabel}\leanok
\uses{def:sigma0}
\end{definition}

\begin{lemma}[Relabelling blocks commutes with rotating within blocks]\label{lem:lex-relabel-comm}
Every block relabelling commutes with $\sigma_0$.
\lean{ConwayO7.Encoder.cycleRelabel_comm, ConwayO7.Encoder.cycleRelabelVal_comm}\leanok
\uses{def:lex-cycle-relabel, def:sigma0}
\end{lemma}

\begin{proof}[Proof sketch]
In the coordinates (block index, position within block) of a nonzero vertex,
$\sigma_0$ increments the position modulo $7$ and fixes the block index,
while a block relabelling permutes the block index and fixes the position;
the two actions touch disjoint coordinates.  Both maps fix the vertex $0$.
\end{proof}

%---------------------------------------------------------------------
\section{The neighbourhood of the fixed vertex}

\begin{fact}[Pair orbits at the fixed vertex]\label{fact:lex-orbit-block}
For $v \neq 0$ the orbit of the pair $\{0,v\}$ under
$\langle\sigma_0\rangle$ depends only on the block of $v$: it equals the
orbit of $\{0, b\}$ where $b$ is the first vertex of $v$'s block.  (Checked
by computation over all $98$ vertices.)
\lean{ConwayO7.Encoder.orbit0_cyc}\leanok
\uses{def:sigma0, lem:orbits}
\end{fact}

\begin{definition}[Block bits]\label{def:lex-block-bit}
For a word $\omega \in \word{693}$ and a block index $c$, the \emph{block
bit} of $c$ is the value of $\omega$ on the orbit of the pair
$\{0, 1+7c\}$, i.e.\ the adjacency bit of the fixed vertex towards
block~$c$.
\lean{ConwayO7.Encoder.cycleBit}\leanok
\uses{lem:orbits, fact:lex-orbit-block}
\end{definition}

\begin{lemma}[Adjacency to a block is constant on the block]\label{lem:lex-adj-block}
In the graph $G_\omega$ associated to a word $\omega$, the fixed vertex $0$
is adjacent to a vertex $v$ if and only if $v \neq 0$ and the block bit of
$v$'s block equals $1$.  In particular the neighbourhood of $0$ is a union
of full blocks.
\lean{ConwayO7.Encoder.adj_zero_iff}\leanok
\uses{prop:orbitfun, def:lex-block-bit, fact:lex-orbit-block}
\end{lemma}

\begin{proof}[Proof sketch]
Adjacency in $G_\omega$ is the value of $\omega$ on the orbit of the pair,
and by Fact~\ref{fact:lex-orbit-block} that orbit depends only on the block
of~$v$.
\end{proof}

\begin{fact}[Block sizes]\label{fact:lex-block-card}
Each of the fourteen blocks contains exactly $7$ vertices.
\lean{ConwayO7.Encoder.cycle_fiber_card}\leanok
\uses{def:sigma0}
\end{fact}

\begin{lemma}[Exactly two neighbour blocks]\label{lem:lex-two-blocks}
If $G_\omega$ is an $\srg(99,14,1,2)$, then exactly two of the fourteen
block bits of $\omega$ equal $1$.
\lean{ConwayO7.Encoder.two_neighbor_cycles}\leanok
\uses{def:srg, lem:lex-adj-block, fact:lex-block-card}
\end{lemma}

\begin{proof}[Proof sketch]
By regularity the fixed vertex has $14$ neighbours; by
Lemma~\ref{lem:lex-adj-block} its neighbourhood is the disjoint union of the
blocks whose bit is $1$, each of size $7$
(Fact~\ref{fact:lex-block-card}).  Counting fiberwise, $14 = 7m$ where $m$
is the number of blocks with bit $1$, so $m = 2$.
\end{proof}

\begin{lemma}[Normalisation of the fixed vertex's neighbourhood]\label{lem:sym1}
Every $\sigma_0$-invariant $\srg(99,14,1,2)$ is isomorphic, via a block
relabelling commuting with $\sigma_0$, to one in which the neighbourhood of
the fixed vertex $0$ is exactly the union of blocks $0$ and $1$: its orbit
word has block bit $1$ exactly at blocks $0$ and $1$.
\lean{ConwayO7.Encoder.sym1_seed, ConwayO7.Encoder.comap_invariant}\leanok
\uses{prop:orbitfun, def:lex-cycle-relabel, lem:lex-relabel-comm, lem:lex-two-blocks, lem:grp-srg-transport}
\end{lemma}

\begin{proof}[Proof sketch]
By Lemma~\ref{lem:lex-two-blocks} the neighbourhood of $0$ is the union of
two blocks $d_1 \neq d_2$.  A product of two transpositions in $S_{14}$
carries $0,1$ to $d_1,d_2$; the induced block relabelling
(Definition~\ref{def:lex-cycle-relabel}) commutes with $\sigma_0$
(Lemma~\ref{lem:lex-relabel-comm}), so the pulled-back graph is again
$\sigma_0$-invariant and hence an orbit-word graph; it is strongly regular
by transport (Lemma~\ref{lem:grp-srg-transport}), and its block bits are
those of $\omega$ permuted so that the two bits equal to $1$ land at blocks
$0$ and~$1$.
\end{proof}

%---------------------------------------------------------------------
\section{The generator action on orbit words}

\begin{definition}[Vertex action of a generator]\label{def:lex-vertex-action}
Each of the $101$ relabelling generators of Definition~\ref{def:relabel} is
given by an explicit map on the vertex set $\{0,\dots,98\}$; this vertex map
induces the generator's action on pair orbits and hence on words.
\lean{ConwayO7.Encoder.permVertex}\leanok
\uses{def:relabel}
\end{definition}

\begin{fact}[Compatibility with the orbit structure]\label{fact:lex-orbit-compat}
For each of the $101$ generators $\rho$ and every pair of distinct vertices
$u,v$, the orbit index assigned by $\rho$ to the orbit of $\{u,v\}$ is the
orbit of $\{\rho(u),\rho(v)\}$: the action on orbit words is induced by the
action on vertices.  (Checked by computation over all $101 \cdot
\binom{99}{2}$ instances.)
\lean{ConwayO7.Encoder.orbit_compat_of_mem, ConwayO7.Encoder.lexPerms_orbit_compat}\leanok
\uses{def:relabel, def:lex-vertex-action, lem:orbits}
\end{fact}

\begin{fact}[Injectivity of the vertex maps]\label{fact:lex-vertex-inj}
The vertex map of each of the $101$ generators is injective, hence a
permutation of the $99$ vertices.  (Checked by computation.)
\lean{ConwayO7.Encoder.vertex_inj_of_mem, ConwayO7.Encoder.lexPerms_vertex_inj}\leanok
\uses{def:relabel, def:lex-vertex-action}
\end{fact}

\begin{lemma}[One generator step preserves strong regularity]\label{lem:lex-step-srg}
If $G_\omega$ is an $\srg(99,14,1,2)$ and $\rho$ is one of the $101$
generators, then $G_{\rho\cdot\omega}$ is again an $\srg(99,14,1,2)$.
\lean{ConwayO7.Encoder.step_srg}\leanok
\uses{prop:orbitfun, fact:lex-orbit-compat, fact:lex-vertex-inj, lem:grp-srg-transport}
\end{lemma}

\begin{proof}[Proof sketch]
By Fact~\ref{fact:lex-vertex-inj} the vertex map of $\rho$ is a bijection of
the vertex set; by Fact~\ref{fact:lex-orbit-compat} it defines a graph
isomorphism $G_{\rho\cdot\omega} \cong G_\omega$.  Conclude by
Lemma~\ref{lem:grp-srg-transport}.
\end{proof}

\begin{lemma}[One generator step preserves the normalisation]\label{lem:lex-step-sym1}
Each of the $101$ generators permutes the fourteen orbits of pairs
$\{0, 1+7i\}$ among themselves and preserves the predicate ``$i < 2$'' on
their indices; consequently, if the word $\omega$ satisfies the
normalisation condition (3) of Definition~\ref{def:phi}, so does
$\rho\cdot\omega$.
\lean{ConwayO7.Encoder.step_sym1, ConwayO7.Encoder.lexPerms_sym1}\leanok
\uses{def:relabel, def:phi, def:lex-vertex-action}
\end{lemma}

\begin{proof}[Proof sketch]
The generators normalise the stabiliser structure at the fixed vertex: each
sends a block either to another block wholesale, and blocks $0,1$ to blocks
$0,1$.  The corresponding permutation of the fourteen distinguished orbit
indices, together with the invariance of ``$i<2$'', is checked by
computation for each generator; the preservation of condition (3) follows by
rewriting each prescribed bit of $\rho\cdot\omega$ as a prescribed bit of
$\omega$.
\end{proof}

\begin{definition}[Reachability]\label{def:lex-reach}
For $\omega \in \word{693}$, let $\mathcal{R}(\omega)$ be the smallest set
of words containing $\omega$ and closed under the action of each of the
$101$ generators: the orbit of $\omega$ under the monoid they generate.
\lean{ConwayO7.Encoder.Reach}\leanok
\uses{def:relabel}
\end{definition}

\begin{lemma}[Reachability preserves the invariants]\label{lem:lex-reach-inv}
If $G_\omega$ is an $\srg(99,14,1,2)$, then so is $G_{\omega'}$ for every
$\omega' \in \mathcal{R}(\omega)$; and if $\omega$ satisfies the
normalisation condition, so does every $\omega' \in \mathcal{R}(\omega)$.
\lean{ConwayO7.Encoder.reach_srg, ConwayO7.Encoder.reach_sym1}\leanok
\uses{def:lex-reach, lem:lex-step-srg, lem:lex-step-sym1}
\end{lemma}

\begin{proof}[Proof sketch]
Induction along the chain of generator steps, using
Lemmas~\ref{lem:lex-step-srg} and~\ref{lem:lex-step-sym1} at each step.
\end{proof}

%---------------------------------------------------------------------
\section{Word values and minimisation}

\begin{definition}[Big-endian value]\label{def:lex-word-val}
The \emph{value} of a word $\omega \in \word{693}$ is
$\mathrm{val}(\omega) = \sum_{i<693} \omega(i)\, 2^{692-i}$, the word read
as a big-endian binary number.
\lean{ConwayO7.Encoder.wordVal}\leanok
\uses{lem:orbits}
\end{definition}

\begin{lemma}[Value order refines lexicographic order]\label{lem:lex-value-order}
Let $u, w \in \word{693}$ with $\mathrm{val}(u) \le \mathrm{val}(w)$.  Then
for every position $i < 693$: if $u$ and $w$ agree at all positions
$j < i$ and $u(i) = 1$, then $w(i) = 1$.  Equivalently, $u \preceq w$ in the
lexicographic order.
\lean{ConwayO7.Encoder.lex_of_val_le}\leanok
\uses{def:lex-word-val}
\end{lemma}

\begin{proof}[Proof sketch]
Suppose $u(i)=1$ but $w(i)=0$ with equal prefixes below $i$.  Split both
values at position $i$: the prefix contributions coincide, the bit of $u$ at
$i$ contributes $2^{692-i}$, and the tail of $w$ is at most
$\sum_{t<692-i} 2^t = 2^{692-i}-1$.  Hence
$\mathrm{val}(u) > \mathrm{val}(w)$, a contradiction.
\end{proof}

\begin{lemma}[Existence of a reachable minimum]\label{lem:lex-min-reach}
For every $\omega$ there is $\omega_{\min} \in \mathcal{R}(\omega)$ with
$\mathrm{val}(\omega_{\min}) \le \mathrm{val}(\omega')$ for all
$\omega' \in \mathcal{R}(\omega)$.
\lean{ConwayO7.Encoder.exists_min_reach}\leanok
\uses{def:lex-reach, def:lex-word-val}
\end{lemma}

\begin{proof}[Proof sketch]
The set of values of reachable words is a nonempty set of natural numbers;
by well-ordering it has a least element, attained by some reachable word.
\end{proof}

%---------------------------------------------------------------------
\section{The comparison chains}

\begin{definition}[Lexicographic comparison of literal words]\label{def:lex-leq}
Fix two sequences $X, Y$ of $693$ literals over the orbit variables — for
the encoding, $X$ lists the orbit variables in position order and $Y$ lists
them re-indexed by a generator's orbit action.  A valuation $a$ satisfies
$X \preceq Y$ if for every $i < 693$: whenever the literal values of $X$ and
$Y$ agree at all positions $j < i$ and the $i$-th literal of $X$ is true,
the $i$-th literal of $Y$ is true.  This is the lexicographic comparison of
the two literal-value words under $a$.
\lean{ConwayO7.Encoder.LexLeq, ConwayO7.Encoder.lexXA, ConwayO7.Encoder.lexYA}\leanok
\uses{def:phi, def:relabel}
\end{definition}

\begin{definition}[The lex-leader invariant]\label{def:lex-inv}
A valuation of the orbit variables satisfies the \emph{lex-leader
invariant} if for each of the $101$ generators $\rho$ it satisfies
$X \preceq Y_\rho$, where $Y_\rho$ is $X$ re-indexed by $\rho$'s orbit
action.  On words this says precisely
$\omega \preceq \rho\cdot\omega$ for every generator $\rho$, i.e.\
condition (4) of Definition~\ref{def:phi}.
\lean{ConwayO7.Encoder.LexInv}\leanok
\uses{def:lex-leq, def:relabel, def:phi}
\end{definition}

\begin{lemma}[Satisfiability of one comparison chain]\label{lem:lex-chain-sat}
The clauses rendering one comparison $X \preceq Y$ use fresh auxiliary
registers.  Every valuation $a$ that satisfies all previously emitted
clauses and the comparison $X \preceq Y$ (on literal values) extends, by
assigning only the fresh registers, to a valuation satisfying all clauses of
the chain as well.
\lean{ConwayO7.Encoder.specAddLexLeq_sat, ConwayO7.Encoder.specAddLexLeq}\leanok
\uses{def:lex-leq, def:phi}
\end{lemma}

\begin{proof}[Proof sketch]
The chain introduces, at each position where the two literals differ
structurally, one register meaning ``all earlier positions agree'', defined
by clauses from its predecessor and the current pair, together with the
implication ($i$-th literal of $X$ true $\Rightarrow$ $i$-th literal of $Y$
true) guarded by the previous register.  Assign each register its intended
truth value — the running equality flag of the literal-value words.  Every
defining clause then holds by construction, and every guarded implication
holds because the comparison $X \preceq Y$ supplies exactly the implication
at each position whose strict prefix agrees.  Induction over the positions
of the chain.
\end{proof}

\begin{proposition}[The comparison section]\label{prop:lex-section-trip}
The full minimality section of the encoding — the $101$ chained comparisons
of condition (4) — preserves satisfiability: any valuation of the orbit
block satisfying the lex-leader invariant extends over the auxiliary
registers of the section to satisfy every clause the section emits.
\lean{ConwayO7.Encoder.lex_section_trip, ConwayO7.Encoder.specAddLexLeq_trip, ConwayO7.Encoder.specLexAgainst_trip}\leanok
\uses{def:lex-inv, lem:lex-chain-sat}
\end{proposition}

\begin{proof}[Proof sketch]
Each of the $101$ chains is handled by Lemma~\ref{lem:lex-chain-sat}, whose
hypothesis is one conjunct of the lex-leader invariant; the extensions
compose because each chain touches only its own fresh registers, leaving the
orbit block — and hence the invariant — unchanged.
\end{proof}

%---------------------------------------------------------------------
\section{Assembly}

\begin{theorem}[Completeness of minimality]\label{thm:lex-complete}
Every $\sigma_0$-invariant $\srg(99,14,1,2)$ admits an isomorphic copy whose
orbit word $\omega$ satisfies conditions (3)--(4) of
Definition~\ref{def:phi}: the normalisation of Lemma~\ref{lem:sym1} and the
lex-leader invariant $\omega \preceq \rho\cdot\omega$ for each of the $101$
generators $\rho$.
\lean{ConwayO7.Encoder.lexLeaderComplete, ConwayO7.Encoder.LexLeaderComplete}\leanok
\uses{lem:sym1, def:lex-reach, lem:lex-reach-inv, lem:lex-min-reach, lem:lex-value-order, def:lex-inv}
\end{theorem}

\begin{proof}[Proof sketch]
Normalise first: by Lemma~\ref{lem:sym1} pass to a word $\omega_0$ that is
strongly regular and normalised.  Among the words reachable from $\omega_0$
(Definition~\ref{def:lex-reach}), choose one of minimal value,
$\omega_{\min}$ (Lemma~\ref{lem:lex-min-reach}); by
Lemma~\ref{lem:lex-reach-inv} it is still strongly regular and normalised.
For each generator $\rho$ the word $\rho\cdot\omega_{\min}$ is again
reachable, so
$\mathrm{val}(\omega_{\min}) \le \mathrm{val}(\rho\cdot\omega_{\min})$, and
by Lemma~\ref{lem:lex-value-order} value order refines lexicographic order:
$\omega_{\min} \preceq \rho\cdot\omega_{\min}$.  Translating the word-level
comparison into the literal-level comparison of
Definition~\ref{def:lex-leq} gives the lex-leader invariant.
\end{proof}

\begin{definition}[The master invariant]\label{def:lex-full-inv}
The \emph{master invariant} of a valuation of the orbit block is the
conjunction of the semantic conditions of all four sections of the
encoding: the degree counts (1), the common-neighbour counts (2), the
normalisation (3), and the lex-leader invariant (4).
\lean{ConwayO7.Encoder.FullInv}\leanok
\uses{def:phi, prop:counters, def:lex-inv}
\end{definition}

\begin{proposition}[Generator satisfiability]\label{prop:lex-generator-sat}
Every valuation of the orbit block satisfying the master invariant extends,
over the auxiliary variables, to a valuation satisfying every clause of the
generated formula $\Phi$.
\lean{ConwayO7.Encoder.mkO7St_sat, ConwayO7.Encoder.fullSpec_trip}\leanok
\uses{def:lex-full-inv, prop:counters, prop:lex-section-trip, def:phi}
\end{proposition}

\begin{proof}[Proof sketch]
The four sections are emitted in sequence, each consuming one conjunct of
the master invariant and preserving the rest (each section's auxiliary
variables are fresh, so the other conjuncts, which depend only on the orbit
block, are untouched).  The counting sections extend by
Proposition~\ref{prop:counters}, the comparison section by
Proposition~\ref{prop:lex-section-trip}; composing the four extension steps
yields a satisfying valuation of the whole formula.  That the abstractly
described pipeline emits exactly the clauses of $\Phi$ is checked by
computation.
\end{proof}

\begin{lemma}[Faithfulness from completeness of minimality]\label{lem:lex-encodes}
Assume completeness of minimality
(Theorem~\ref{thm:lex-complete}).  If a $\sigma_0$-invariant
$\srg(99,14,1,2)$ exists, then the distributed formula is satisfiable.
\lean{ConwayO7.Encoder.encodes_of_lexLeader}\leanok
\uses{prop:orbitfun, lem:grp-srg-transport, def:lex-full-inv, prop:lex-generator-sat, fact:bytes}
\end{lemma}

\begin{proof}[Proof sketch]
An invariant graph descends to an orbit word $\omega$
(Proposition~\ref{prop:orbitfun}) whose associated graph is strongly
regular by transport (Lemma~\ref{lem:grp-srg-transport}).  Completeness
of minimality replaces $\omega$ by a companion word satisfying strong
regularity, the normalisation, and the lex-leader invariant; the arithmetic
conditions (1)--(2) hold for the orbit word of any strongly regular graph,
so the companion satisfies the master invariant.
Proposition~\ref{prop:lex-generator-sat} extends it to a satisfying
valuation of the generated formula, which equals the distributed file by the
byte identity (Fact~\ref{fact:bytes}).
\end{proof}

\begin{theorem}[Completeness of the encoding]\label{thm:encoding-complete}
If a $\sigma_0$-invariant $\srg(99,14,1,2)$ exists, then $\Phi$ is
satisfiable.
\lean{ConwayO7.Encoder.encodes_o7}\leanok
\uses{thm:lex-complete, lem:lex-encodes}
\end{theorem}

\begin{proof}[Proof sketch]
Combine Theorem~\ref{thm:lex-complete} with Lemma~\ref{lem:lex-encodes}.
\end{proof}

%=====================================================================
